\section{Psychology of Mystical Experience}

\begin{quotex}
We can indeed know that the bosom of the Father is our own basis and origin, and that our life and being find therein their life and principle. From this basis that is proper to us---that is, from the bosom of the Father, and from all that lives in him---shines forth an eternal brightness, the generation of the Son. And in this brightness, which is the Son, God sees himself openly, with all that lives in him. All those who, above their created being, are raised to a contemplative life, are one with this divine brightness. They are that brightness itself, and they see, feel and discover, under this divine light, that according to their ideal or uncreated being, they are themselves this abyss of simplicity, the brightness whereof shines without means in divine modes. Thus the contemplatives attain their eternal exemplar, after whose image they have been created, and they contemplate God in all things, without distinction, by a simple gaze, in the divine brightness.
\flright{\textsc{John of Ruysbroeck}}
\end{quotex}


The concluding chapter of \emph{Studies in the Psychology of the Mystics} by \textbf{Joseph Marechal}, S.J. deals with how one's world conception colours his mystical experience. As a methodological principle, Fr. Marechal is not concerned with the truth of the metaphysical systems, but rather with its psychological effects. The obvious problem with his study is how to define ``mystical experience'', since Christian ecstasy and yogic Samadhi are often included as such experiences along with trances induced by hallucinogens, ritual dance, etc. This, then, is his definition of \textbf{mystical experience}:

\begin{quotex}
Mystical experience is a religious experience which is esteemed as superior to the normal: more direct, more intimate, or more rare.
\end{quotex}


The corollary is that there are three fundamental elements:

\begin{enumerate}


\item A religious doctrine---rational or revealed---which is a metempirical doctrine, relative to the Absolute.

\item Psychological facts of actual experience that are relatively rare or exceptional, and susceptible of a religious interpretation.

\item A synthesis of doctrine and psychology which interprets the psychological facts as a function of the doctrine.
\end{enumerate}


Not unlike Rene Guenon, Fr. Marechal accepts the ideal of a unique metaphysic, despite the historical existence of a variety of metaphysical systems. Given the three elements, the highest and purest expressions of mystical experience will be associated with an intellectual understanding of the ideal metaphysic. Fr. Marechal describes several metaphysical systems and the mystical experiences arising from them.

\paragraph{Animism}


Inferior systems have not yet achieved a unity of thought from the plurality of sense impressions. An example is animism which postulates metempirical ``souls'' behind the world of outward experience. The medicine man or sorcerer will associate strange subjective states with these hidden entities. This includes experiences, among others, such as:

\begin{itemize}


\item Dreams with a sacred content

\item Narcotic-induced delirium

\item Trances resulting from extreme physical ascesis

\item Hysterical ecstasies brought on by music or scents
\end{itemize}


Although such experiences are valued very highly and are still sought after by a substantial number of spiritual seekers, Fr. Marechal considers them not worthy of an extended treatment.

\paragraph{Metaphysical Awakening}


\hfill\break

Fr. Marechal points out the close connection of metaphysics with religion in general and with mysticism in particular. He use the famous panel from the Sistine Chapel showing Adam's attraction to finger of God pointing at him as the artistic rendition of the awakening of human intelligence to the notion of an Absolute distinct from himself. It shows:

\begin{itemize}


\item Man's consciousness of his own ego, at once infinite and limited

\item His perception of the universe

\item His thought seeking the Absolute
\end{itemize}


This is the Great Triad: the ego, the universe, and the Absolute. The relationships of these terms are the sources of two fundamental things:

\begin{enumerate}


\item The source of his potentialities of action in the physical order

\item The law governing his attitude and the key that unlocks his destiny in the moral order
\end{enumerate}


Different metaphysical systems define those relationships in different ways. Although Fr. Marechal does not mention this, it is obvious that, as a practical matter, two sciences are a great help in developing those potentialities: esoteric cosmology, which deals with the universe, and esoteric psychology, which deals with man's inner states of consciousness.

\paragraph{Dualism and Pessimism}


Dualist systems see the universe as the expression of the conflict of two eternal principles: e.g., Being or the Good vs non-being or Evil. Such viewpoints lead to destructive asceticism and negative mysticism. That is, one's destiny is determined more by the attempted renunciation of the ``evil spirit'' rather than by a path of positive and uplifting stages. The world is experienced a being limited by the unreal or not-being. The universe is seen as evil and illusory.

Marechal refers to \textbf{Arthur Schopenhauer}`s philosophical system, although not strictly a dualism, but a pessimism. The will to live creates a universe of ``my'' representations, which is nevertheless evil and meaningless. Only the renunciation of that will can be the solution. Again, beyond the negation of phenomena, there is nothing positive. In Manichaean, and even Buddhist mysticism, object, actions, etc., are all limitations; that is the sole Evil. Such mystical practices focus on overcoming the limitations of the World of Becoming, without consideration of what lies beyond in the World of Being.

\paragraph{Pantheistic Monism}


This is a more advanced of sophisticated metaphysical system, what we may call the God of the philosophers. There is no opposition to Being; rather primacy belongs to the Absolute, which is beyond dualism. The earliest such systems derive from the Upanishads or Vedanta. The essential basis of this philosophy is the identity of Brahman and the Atman, which is the sole reality. The apparent multiplicity of things and souls is illusory. Hence, man's destiny is extirpate in his soul the illusory multiplicity of objects and acts to end up within the Atman. The Atman is known, but not as an object. Hence, the only goal is to ``be'' the Atman.

Comparable Western systems never quite go that far, since the idea of an ascent to God via created things is not denied. The influence of Plotinus affects not only the West, but Sufism and even Indian philosophers. God, or the One, eternally creates the world by casting out his rays like a sun to the ``very confines of non-being''. Then the things, which are divine fragments, have a desire to return to the unity of the absolute Good.

The human soul is between the pure Ideas and matter; it is attracted towards the Center by Love. If love predominates in the soul, it begins to be concentrated. Firstly, it apprehends intelligible beauty through the senses. Then by contemplation, it is purified and unified: ``to contemplate is to become what is contemplated.'' From sensible beauty, the soul passes to the lower intellect of concepts, or psyche. Then it becomes the higher Intelligence, or nous.

Still the drive of Eros leads it to seek the perfect unity of the absolute Good. Thus it contemplated Being as the summit of the intelligible world. Yet, there is still the duality of Essence and Existence. Ultimately, though Love, even that duality is overcome in the One.

Later metaphysicians such as \textbf{Giordano Bruno}, \textbf{Benedict Spinoza}, or \textbf{Johann Fichte} are in this stream. For example, in the negative direction, Spinoza defines man's moral end as the freedom gained over inadequate ideas and the passions. Yet, there is the positive direction of the ``intellectual love of God''.

Fichte, from whom Evola borrowed much, discovers the pure I, or spirit, opposing to itself the not-I, or universe, so that it may progressively know itself and conquer itself by gaining self-mastery. The human intelligence is raised to the dignity of an Absolute, i.e., a ``mysticism of becoming God.''

The point is that ``metaphysics can invest the ordinary operations of our understanding with a mystical significance.'' However, Marechal identifies the mysticism that such systems inspire as purely ``natural'', achieved through thought, and lacking any notion of ``grace''. In short, the distinction between nature and supernature is obscured.

\paragraph{Monotheism and Supernatural Mysticism}


Finally, Marechal address strict monotheism as the answer to the question of the relationships between the Ego, the universe, and the Absolute. God, then, is strictly transcendent without any common measure with that which is not himself. He possesses the fullness of Being which excludes not-Being.

The universe then is a ``becoming'' proceeding from God and arranged between the limits of pure Being and pure Not-Being. Although external to God, it is moved and directed by divine action. Finite things do not create themselves as in in monistic systems; yet they are not evil or nothingness as in the pessimist dualist systems. They are free creations of love and possess value measured by their degree of participation in the perfection of the divine Being.

Hence, the created intelligences and wills will tend to reproduce the divine Ideal and try to make progress towards it. Hence, the mysticism from this view will have many analogies with Neoplatonism. Creation is grafted onto the closed cycle of the operations of Divine creation as an epicycle. At the beginning of the epicycle, the divine actions descends into the innermost elements of things, giving them their ``nature''.

Next the flood ascends back to God, as every material beings ``tends to the perfection of its species''. The entire physical world tends toward vital unity in stages:

\begin{itemize}


\item The vital unity, through its comprehensive interiority, foretells and prepares for consciousness

\item Sensibility (or sense experience) reflects the unconscious world of matter and brings it to the threshold of the idea

\item Then, the intelligence recognizes the intelligible in the data of sense

\item Finally, it discovers in its ``becoming'' no other end except God himself. This closes the epicycle.
\end{itemize}


The return of things to their first principle, i.e., the ascent of the intelligence to God, is a mystical phase. All becoming has a \emph{law}, and in that law, is an \emph{end}. Marechal then describes that law as it is fostered in the depth of the human soul. There are two aspects: the order of intelligence and the order of the will.

The Intellect tends toward the assimilation of Being as known, or, Absolute Truth. The Will tends to the possession of Being as absolute Good. The two tendencies converge on the direct vision of God. However, insofar as the soul is restricted to the plane of creation, there must also be an initiative from God that expands the intelligence and reveals himself.

This initiative from above leads to higher states than what can be achieved through the natural mysticism of monism, as it opens up the horizons of grace and supernature. It does not reject the natural metaphysical systems, but builds on them. Paradoxically, it leads to unity with the divine while retaining the individuality of the soul. That is the highest achievement of the Western Tradition.

\paragraph{Summary}


Marechal summarizes the three principal types of mystical theories:

\begin{enumerate}


\item \textbf{Negative mysticism}. A mysticism of simple liberation, the result of dualistic cosmologies and philosophic pessimism.

\item \textbf{Positive mysticism}. A mysticism of divine becoming immanent in the soul.

\item \textbf{Theistic mysticism}. The mysticism of objective striving toward God by means of knowledge and love.
\end{enumerate}


Clearly, a superior metaphysic will lead to deeper mystical experiences.


\flrightit{Posted on 2016-01-27 by Cologero}

\begin{center}* * *\end{center}

\begin{footnotesize}
\begin{sffamily}
\texttt{Mercurius on 2016-01-28 at 22:59 said:}

Another excellent piece Cologero.

Concerning ``pantheistic monism'', the West, and the manner in which Western ``monotheism and supernatural mysticism'' includes the former, while perfecting and it and going beyond, as contradictory it might seem, the significance of a pantheistic materialist monism ought factor in as well. Vis a vis, what the Stoa brought to the table.

But, how can THAT be of any use to us, we might ask--as the strict form excludes any transcendence (apparently), contradicts Platonism, is formally anti-Idealistic, and insists on an endless identical ``recur''.

What is brilliant, is that Paul (seems to have) taken Stoa materialistic pantheism, endorsed certain basic implications in it, while correcting and adjusting for errors; and consequently, the Western Tradition's ``dna'', counterbalances and adjusts for pure idealist errors.

In 1 Corinthians, Paul wrote: ``Now you are the body of Christ, and each one of you is a part of it''. Here, a sort of material ``pantheism'' is expressed, while removed from a mere naturalism.

\hfill

\texttt{Mark Citadel on 2016-02-03 at 12:25 said:}

Well this was a joy to read as always. A couple of queries:

``This is the Great Triad: the ego, the universe, and the Absolute.''

An odd question, but do angelic beings find themselves in this relationship as egos, a mechanic part of the universe for man, or indeed a functioning branch of the absolute (i.e -- the Divine Realm). I am tempted to say they are egos as well, since we know that angels had (or still have) free will and moral agency.

Also, on dualism, in Scripture we are given a dichotomy between choosing `the World' or choosing God. I am curious as to how we can overcome the seemingly dualistic nature of this. Is it that we must see `the World' not as an expression of evil, but simple `becoming' as opposed to the superior `being'?

``Hence, the created intelligences and wills will tend to reproduce the divine Ideal and try to make progress towards it.''

This was what I had theorized when pontificating on the ever-present nature of `Tradition' that in fact installations of monarchy and patriarchy are reflections of the Divine Realm, aspired to by all peoples, hence their ubiquity. However, right now we seem to be progressing AWAY from that ideal to some horrific nadir. In your four stages... ``- The vital unity, through its comprehensive inferiority, foretells and prepares for consciousness\\ -- Sensibility (or sense experience) reflects the unconscious world of matter and brings it to the threshold of the idea\\ -- Then, the intelligence recognizes the intelligible in the data of sense\\ -- Finally, it discovers in its ``becoming'' no other end except God himself. This closes the epicycle.''

Where would the current state of man in the Kali Yuga rest? Or, at the bottom of the cycle, have we not even a foot on ladder rung \#1!

\hfill

\texttt{Cologero on 2016-02-04 at 22:39 said:}

@Mark, not sure what you are driving at. Angelic beings have a self, though not like man's ego. Guenon explained that angels could be considered as higher states of being. That is, they are beyond sense impressions and discursive knowledge. They have the beatific vision and do not need to seek it.

I think your initial impression is correct that the ``world'' refers to the lack of a sense of transcendence. As such, ``the World'' is not a metaphysical concept as the world of becoming vis-a-vis the world of Being. Otherwise, your best bet would be idealist monism, just contemplating the ideal world of Being.

It is not a question of ``we'', it starts with a few. Many are at stage 4. The Kali Yuga will end whenever you want it to.

\hfill

\texttt{Miltie on 2016-02-06 at 21:39 said:}

Changing the ``a'' to an ``o'' gives a very very different translation in Greek.

\hfill

\texttt{Cologero on 2016-02-06 at 22:38 said:}

Uncle Miltie, FWIW, I've always been told I have a nice butt. I think I'll change my name to kalokolo.

\end{sffamily}
\end{footnotesize}

